\begin{Theorem}
\label{thm:compact-perturbation-fredholm-index-invariance}
设 $A\in\mathcal{L}(X,Y)$ 为半 Fredholm 算子, $k\in\mathcal{L}(X,Y)$ 为紧算子. 那么 $A+k$ 为半 Fredholm 算子, 且 $\operatorname{ind}A+k=\operatorname{ind}A$.
\end{Theorem}

\begin{Proof}
先考虑 $Y=X$ 且 $A=I$ 的特殊情形. 需要证明 $I+k$ 为指标为零的 Fredholm 算子. 设 $x_n\in\mathcal{N}(I+k)$, 且 $\|x_n\|\leq1$. 抽取子序列后, 可假设 $kx_n$ 收敛, 又由于 $x_n+kx_n=0$, 得到 $x_n$ 收敛. 因而 $\mathcal{N}(I+k)$ 中的单位球是紧的, 从而 $\mathcal{N}(I+k)$ 有限维.

接着证明 $\mathcal{R}(I+k)$ 是闭的. 设 $V$ 为 $\mathcal{N}(I+k)$ 的补空间, 那么 $I+k$ 为从 $V$ 到 $\mathcal{R}(I+k)$ 的双射. 若证明其逆有界, 则 $\mathcal{R}(I+k)$ 与 $V$ 同构, 因而完备. 现在假设其逆无界, 即存在序列 $x_n\in V$, 使得 $\|x_n\|=1$, 但 $x_n+kx_n\to0$. 那么可以再次抽取子序列, 假设 $kx_n$ 收敛, 记其极限为 $y$. 于是 $x_n\to-y$, 且 $y+ky=0$, 即 $y\in\mathcal{N}(I+k)$. 但这产生矛盾, 因为另一方面 $y\in V$ 且 $\|y\|=1$.

这样就证明了 $I+k$ 为半 Fredholm 算子. 考虑族 $\{I+tk\mid t\in[0,1]\}$, 即知指标为零.

对于一般情形, 仍假设 $A$ 为 Fredholm 算子, 或 $X$ 和 $Y$ 为 Hilbert 空间. 设 $X=V\oplus\mathcal{N}(A)$, $Y=W\oplus\mathcal{R}(A)$. 如下定义 $S\in\mathcal{L}(Y,X)$: 若 $y\in\mathcal{R}(A)$, 则 $Sy$ 为满足 $Ax=y$ 的唯一 $x\in V$; 若 $y\in W$, 则 $Sy=0$. 若 $W$ 有限维, 则 $AS$ 为恒等算子的有限秩扰动, 即 $AS-I$ 紧. 因而 $(A+k)S-I$ 也紧, 即 $(A+k)S$ 为指标为零的 Fredholm 算子. 所以 $\mathcal{R}((A+k)S)$ 具有有限余维, 从而 $\mathcal{R}(A+k)$ 也具有有限余维.

另一方面, 若 $\mathcal{N}(A)$ 有限维, 则 $SA-I$ 为有限秩算子, $S(A+k)$ 为指标为零的 Fredholm 算子. 因而 $S(A+k)$ 的零空间有限维, 更不用说 $A+k$ 的零空间也有限维. 需要证明 $A+k$ 的值域是闭的. 设 $U$ 为 $\mathcal{N}(A)$ 和 $\mathcal{N}(A+k)$ 的线性张成空间在 $X$ 中的补空间. 将证明 $A+k|_U$ 有有界逆. 这说明 $(A+k)(U)$ 完备, 又由于该空间在 $\mathcal{R}(A+k)$ 中具有有限余维, 所以 $\mathcal{R}(A+k)$ 闭. 现在假设 $x_n\in U$, 且 $\|x_n\|=1$, $(A+k)x_n\to0$. 抽取子序列后, 可假设 $kx_n$ 收敛, 因而 $Ax_n$ 收敛. 但 $A|_U$ 有有界逆, 所以 $x_n$ 收敛, 记极限为 $y$. 如前所述, 这产生矛盾, 因为 $(A+k)y=0$, 但另一方面 $y\in U$ 且 $\|y\|=1$.

因而已经证明 $A+k$ 为半 Fredholm 算子, 再次考虑族 $\{A+tk\mid t\in[0,1]\}$, 就得到关于指标的结论.
\end{Proof}

现在设 $A$ 为 Hilbert 空间 $H$ 中稠密定义的闭对称算子. 若 $\lambda\in\mathbb{C}\backslash\mathbb{R}$, 则 $A-\lambda I$ 有有界逆; 因而它是单射且具有闭值域. 由 Corollary~\ref{cor:semi-fredholm-open-set-constant-index}, $A-\lambda I$ 的亏数在上半平面和下半平面上分别为常数. 因而, 当 $\lambda$ 位于上半平面时, $A-\lambda I$ 的 Fredholm 指标等于 $A$ 的亏指数 $\gamma^+$ 的负值, 下半平面中的 Fredholm 指标则等于 $-\gamma^-$ (参见 Definition~\ref{def:operator-deficiency-indices}). 由 Theorem~\ref{thm:essential-self-adjointness-zero-deficiency}, $A$ 自伴随, 当且仅当两个亏指数均为零.

\Needspace{5\baselineskip}\textbf{Problems}

\begin{enumerate}[label=\thesection.\arabic*.,ref=\thesection.\arabic*,start=54,series=book-problems]

\item\label{ex:closed-plus-finite-dimensional-subspace} 设 $V$ 为 $X$ 的闭子空间, $W$ 为有限维子空间. 证明 $V+W$ 在 $X$ 中是闭的.

\item\label{ex:positive-potential-dirichlet-classical-solvability} 设 $c\in C([0,1])$, 且 $c(x)\geq0$. 证明, 对每个给定的 $f\in C([0,1])$, 满足边界条件 $y(0)=y(1)=0$ 的方程 $y\prime\prime-c(x)y=f(x)$ 在 $C^2([0,1])$ 中有解.

\item\label{ex:half-line-differentiation-deficiency-indices} 设 $H=L^2(0,\infty)$, $Au=iu\prime$, 定义域为 $H_0^1(0,\infty)$. 求 $A$ 及其伴随的亏指数.

\item\label{ex:real-resolvent-implies-self-adjointness} 设 $A$ 为闭的, 稠密定义的对称算子, 且 $A$ 的预解集至少包含一个实数. 证明 $A$ 自伴随. 若 $A$ 还正定, 证明负实轴属于 $A$ 的预解集.

\end{enumerate}
